% ============================================================
% CHƯƠNG 2: CƠ SỞ LÝ THUYẾT
% ============================================================
\chapter{Cơ sở lý thuyết}
\label{chap:theory}

\section{Tổng quan về nhận dạng cảm xúc giọng nói}
\label{sec:ser_overview}

Nhận dạng cảm xúc giọng nói (Speech Emotion Recognition -- SER) là bài toán phân loại trạng thái cảm xúc của người nói dựa trên tín hiệu âm thanh \cite{el2011, schuller2018}. Quá trình này bao gồm ba bước chính: (1) thu thập và tiền xử lý tín hiệu âm thanh, (2) trích xuất đặc trưng (feature extraction), và (3) phân loại cảm xúc bằng mô hình học máy hoặc học sâu.

Các hệ thống SER hiện đại thường sử dụng đặc trưng dạng phổ (spectral features) kết hợp với mạng nơ-ron sâu (DNN/CNN) để đạt hiệu suất cao. Cách tiếp cận phổ biến nhất là chuyển đổi tín hiệu âm thanh thành dạng biểu diễn hình ảnh (spectrogram) rồi áp dụng các kỹ thuật thị giác máy tính (computer vision) để phân loại.

\section{Xử lý tín hiệu âm thanh số}
\label{sec:signal_processing}

\subsection{Tín hiệu âm thanh và lấy mẫu}

Âm thanh là sóng cơ học lan truyền trong không khí. Khi số hóa, tín hiệu liên tục được lấy mẫu (sampling) tại tần số cố định. Theo định lý Nyquist, tần số lấy mẫu phải gấp đôi tần số cao nhất cần ghi nhận. Trong đồ án này, tất cả các tệp âm thanh được chuẩn hóa về tần số lấy mẫu $f_s = 16\,\text{kHz}$.

\subsection{Bộ lọc tiền nhấn (Pre-emphasis Filter)}

Bộ lọc tiền nhấn là bộ lọc thông cao bậc nhất. Bộ lọc này có tác dụng cân bằng phổ tần số bằng cách khuếch đại dải tần số cao vốn thường bị suy giảm nhanh khi truyền giọng nói trong không khí. Phép biến đổi này giúp nhấn mạnh các phụ âm cùng các biểu cảm âm học tinh tế, đồng thời cải thiện đáng kể tỉ lệ tín hiệu trên nhiễu (SNR) ở dải tần số cao. Công thức toán học biểu diễn bộ lọc tiền nhấn được xác định như sau:
\begin{equation}
    y[n] = x[n] - \alpha \cdot x[n-1]
    \label{eq:preemphasis}
\end{equation}
trong đó $n$ là chỉ số mẫu rời rạc và $\alpha = 0.97$ là hệ số tiền nhấn, với giá trị $\alpha$ gần 1 sẽ giúp nhấn mạnh dải tần số cao một cách rõ rệt.

\subsection{Biến đổi Fourier ngắn hạn (STFT)}

Biến đổi Fourier ngắn hạn (Short-Time Fourier Transform -- STFT) là phương pháp phân tích tín hiệu đồng thời trên cả miền thời gian và miền tần số. Tín hiệu được chia thành các khung (frame) ngắn, mỗi khung được nhân với hàm cửa sổ rồi thực hiện biến đổi Fourier rời rạc (DFT):

\begin{equation}
    X(m, k) = \sum_{n=0}^{N-1} x_{\text{padded}}(n + mH) \cdot e^{-j2\pi kn/N}
    \label{eq:stft}
\end{equation}
Trong đó, khung tín hiệu gốc có độ dài $L$ được nhân với hàm cửa sổ Hamming $w(n)$ và thực hiện đệm thêm $N - L$ giá trị $0$ (zero-padding) để tạo thành chuỗi tín hiệu đệm $x_{\text{padded}}$ có độ dài $N$ trước khi tính toán DFT:
\begin{equation}
    x_{\text{padded}}(n + mH) = \begin{cases} 
      x(n + mH) \cdot w(n) & \text{với } 0 \le n < L \\
      0 & \text{với } L \le n < N 
   \end{cases}
\end{equation}
Trong các công thức trên, các ký hiệu cụ thể được định nghĩa như sau: $m$ là chỉ số khung thời gian; $k$ là chỉ số tần số ($0 \le k < N$); $L$ biểu thị độ dài khung tín hiệu gốc (thiết lập ở mức $640$ mẫu, tương đương $40\,\text{ms}$); $N$ là độ dài cửa sổ FFT ($800$ điểm); $H$ là bước trượt (hop length, thiết lập bằng $160$ mẫu, tương đương $10\,\text{ms}$); và $w(n)$ là hàm cửa sổ Hamming. Các tham số STFT cụ thể được tổng hợp chi tiết trong Bảng~\ref{tab:stft_params}.

\begin{table}[H]
\centering
\caption{Cấu hình tham số STFT}
\label{tab:stft_params}
\begin{tabular}{lll}
\toprule
\textbf{Tham số} & \textbf{Giá trị} & \textbf{Mô tả} \\
\midrule
Window type & Hamming & Hàm cửa sổ \\
Window length & 40 ms (640 mẫu) & Độ dài mỗi khung \\
Hop length & 10 ms (160 mẫu) & Bước trượt giữa các khung \\
N\_FFT & 800 điểm & Số điểm FFT \\
N\_freq & 200 & Số băng tần số giữ lại \\
\bottomrule
\end{tabular}
\end{table}

\subsection{Log-Spectrogram}

Phổ biên độ thu được từ STFT được chuyển sang thang đo logarit (decibel) nhằm mô phỏng sát nhất cách tai người cảm nhận cường độ âm lượng:
\begin{equation}
    S_{\text{dB}}(m, k) = 20 \cdot \log_{10}\left(|X(m, k)|\right)
    \label{eq:logspec}
\end{equation}
Hệ thống chỉ tiến hành giữ lại 200 băng tần số thấp nhất (trên tổng số 401 băng tần của STFT) vì dải tần số dưới $4\,\text{kHz}$ chứa hầu hết các thông tin ngữ điệu cảm xúc quan trọng của con người. Kết quả đầu ra là một ma trận đặc trưng Log-Spectrogram kích thước $200 \times T$ cho mỗi câu nói, với $T$ đại diện cho số lượng khung thời gian phụ thuộc vào thời lượng câu phát âm gốc.


\section{Mạng nơ-ron tích chập (CNN)}
\label{sec:cnn}

\subsection{Kiến trúc tổng quát và Lịch sử phát triển}

Mạng nơ-ron tích chập (Convolutional Neural Network -- CNN) là lớp kiến trúc học sâu được thiết kế để xử lý dữ liệu cấu trúc dạng lưới như hình ảnh hoặc phổ đồ âm thanh số. Sơ đồ kiến trúc tổng quan của một mạng CNN tiêu chuẩn được mô tả trong Hình~\ref{fig:cnn_general_arch}.

\begin{figure}[H]
    \centering
    \includegraphics[width=0.85\textwidth]{cnn_general_arch.png}
    \caption{Sơ đồ kiến trúc tổng quan của mạng nơ-ron tích chập (CNN)}
    \label{fig:cnn_general_arch}
\end{figure}

Kiến trúc CNN tiêu chuẩn bao gồm ba cấu phần cơ bản xếp chồng tuần tự. Trước tiên là các lớp tích chập (Convolutional Layers) áp dụng các bộ lọc (kernel) trượt trên ảnh đầu vào để trích xuất đặc trưng cục bộ. Tiếp theo là các lớp gộp (Pooling Layers) có chức năng giảm kích thước không gian bản đồ đặc trưng bằng phép lấy cực đại (Max Pooling) hoặc trung bình (Average Pooling) nhằm giảm chi phí tính toán và tăng tính bất biến đối với các phép dịch chuyển nhỏ. Cuối cùng là các lớp kết nối đầy đủ (Fully Connected Layers) thực hiện kết hợp các đặc trưng đã trích xuất để phân loại cảm xúc.

Lịch sử phát triển của mạng CNN bắt nguồn từ mô hình Neocognitron đề xuất bởi Kunihiko Fukushima vào năm 1980, giới thiệu các khái niệm sơ khai về lớp trích xuất đặc trưng (S-cells) và lớp bất biến dịch chuyển (C-cells) dựa trên cấu trúc vỏ não thị giác của động vật. Đến năm 1998, Yann LeCun và cộng sự đã hiện thực hóa mô hình này thành kiến trúc LeNet-5 để nhận dạng chữ số viết tay trên séc ngân hàng, đặt nền móng chính thức cho cơ chế tích chập và gộp đặc trưng chia sẻ trọng số. Mặc dù đạt thành công lớn, CNN thời kỳ này bị giới hạn hiệu năng do thiếu tài nguyên tính toán và các kỹ thuật tối ưu hóa hiện đại. 

Bước ngoặt lịch sử diễn ra vào năm 2012 khi kiến trúc AlexNet \cite{krizhevsky2012} tận dụng sức mạnh tính toán song song của GPU để huấn luyện mạng CNN sâu hơn trên tập dữ liệu ImageNet khổng lồ. Kể từ đó, sự phát triển của CNN tăng tốc vượt bậc với sự ra đời của cơ chế kết nối tắt trong ResNet (2016) giúp khắc phục vấn đề triệt tiêu gradient khi mạng đạt độ sâu hàng trăm lớp, cơ chế liên kết dày đặc trong DenseNet (2017) giúp tăng cường tái sử dụng đặc trưng, kiến trúc tích chập tách biệt chiều sâu trong MobileNet (2018) tối ưu cho thiết bị di động, và phương pháp tỷ lệ hóa hỗn hợp trong EfficientNet (2019) giúp cân bằng tối ưu tài nguyên tính toán. Trong những năm gần đây, xu hướng nghiên cứu đang tiếp tục phát triển tới các mô hình kết hợp lai giữa CNN và Vision Transformer (ViT) nhằm khai thác đồng thời khả năng biểu diễn cục bộ của tích chập và khả năng thiết lập mối quan hệ toàn cục của cơ chế attention.

\subsection{AlexNet}

Kiến trúc AlexNet được Krizhevsky và cộng sự đề xuất vào năm 2012 \cite{krizhevsky2012} là một bước ngoặt vĩ đại của học sâu, bao gồm 5 lớp tích chập xếp chồng kết hợp với 3 lớp kết nối đầy đủ. Điểm nổi bật của AlexNet là việc áp dụng hàm kích hoạt ReLU để giải quyết vấn đề triệt tiêu gradient và kỹ thuật Dropout để chống quá khớp trên các lớp kết nối đầy đủ có số lượng tham số lớn. Sơ đồ kiến trúc tổng thể của AlexNet được minh họa chi tiết trong Hình~\ref{fig:alexnet_arch}.

\begin{figure}[H]
    \centering
    \includegraphics[width=0.85\textwidth]{alexnet_arch.png}
    \caption{Sơ đồ kiến trúc mạng nơ-ron tích chập AlexNet}
    \label{fig:alexnet_arch}
\end{figure}

Phân tích sơ đồ Hình~\ref{fig:alexnet_arch} cho thấy mô hình xử lý hình ảnh qua các khối tích chập lớn ở giai đoạn đầu để nắm bắt các đặc trưng hình học thô (như bộ lọc kích thước $11\times11$ ở lớp Conv1 và $5\times5$ ở lớp Conv2), sau đó giảm dần kích thước bộ lọc xuống $3\times3$ ở các lớp Conv3, Conv4 và Conv5 nhằm tinh chỉnh các đặc trưng ngữ nghĩa mức cao. Các lớp gộp cực đại (Max Pooling) được chèn xen kẽ để giảm độ phân giải không gian của bản đồ đặc trưng. Điểm yếu lớn của kiến trúc này nằm ở ba lớp kết nối đầy đủ (FC6, FC7, FC8) phía cuối, nơi chiếm tới hơn 90\% tổng số tham số của mô hình (tạo ra khoảng 57 triệu tham số), khiến mô hình cực kỳ nhạy cảm với hiện tượng quá khớp nếu dữ liệu huấn luyện không đủ lớn.

Trong nghiên cứu này, AlexNet đóng vai trò là mô hình baseline đối chứng với hai phiên bản triển khai cụ thể bao gồm AlexNet Finetuning và AlexNet GAP. Phiên bản AlexNet Finetuning giữ nguyên kiến trúc nguyên bản và chỉ thực hiện điều chỉnh số chiều đầu ra của lớp Fully Connected cuối cùng tương ứng với 4 lớp cảm xúc của bộ dữ liệu IEMOCAP, dẫn đến tổng số tham số lên tới 57,02 triệu. Cần chú ý rằng, mặc dù đầu vào phổ đồ trong đồ án là $256 \times 256$ pixels (vốn làm đầu ra của lớp tích chập Conv5 đạt kích thước $7 \times 7$ thay vì $6 \times 6$ như trong thiết kế ImageNet gốc), việc chuyển giao tri thức từ mô hình pre-trained vẫn thực hiện thành công nhờ việc tích hợp lớp gộp trung bình thích ứng (\texttt{nn.AdaptiveAvgPool2d((6, 6))}) chèn giữa phần trích xuất đặc trưng và bộ phân loại kết nối đầy đủ, giúp đồng bộ hóa kích thước tensor đầu vào cho lớp FC6 mà không làm đứt gãy hoặc yêu cầu tái cấu trúc các trọng số đã được học trước. 

Ngược lại, phiên bản cải tiến AlexNet GAP loại bỏ hoàn toàn các lớp Fully Connected nặng nề và thay thế bằng lớp gộp trung bình toàn cục (Global Average Pooling - GAP) kết hợp tích chập $1\times1$, giúp giảm số lượng tham số huấn luyện xuống chỉ còn 2,47 triệu (giảm 96\% so với nguyên bản) và hạn chế đáng kể hiện tượng quá khớp trên tập dữ liệu nhỏ của bài toán SER. Redesign này giúp cải thiện đáng kể khả năng tổng quát hóa trên dữ liệu kiểm thử.


\section{Transfer Learning}
\label{sec:transfer_learning}

Transfer Learning là kỹ thuật tận dụng trọng số đã được huấn luyện trước từ một tác vụ lớn (thường là ImageNet \cite{deng2009}) để chuyển giao tri thức sang một tác vụ mới có lượng dữ liệu hạn chế. Trong đồ án này, các mô hình CNN tiên tiến được khởi tạo bằng trọng số pre-trained ImageNet, sau đó thực hiện fine-tune trực tiếp trên phổ đồ Log-Spectrogram của dữ liệu IEMOCAP để nhận dạng cảm xúc.

\subsection{ResNet-18}

Mạng dư thừa (Residual Network - ResNet) do He và cộng sự phát triển năm 2016 \cite{he2016} đã giải quyết triệt để hiện tượng suy thoái hiệu năng mạng (degradation problem) khi tăng độ sâu. Cơ chế cốt lõi của ResNet là việc tích hợp các kết nối tắt (skip connection hoặc residual connection) cho phép thông tin đầu vào $\mathbf{x}$ bỏ qua một hoặc nhiều lớp tích chập và được cộng trực tiếp vào đầu ra của khối. Công thức toán học biểu diễn khối dư thừa cơ bản được xác định như sau:
\begin{equation}
    \mathbf{y} = \mathcal{F}(\mathbf{x}, \{W_i\}) + \mathbf{x}
    \label{eq:residual}
\end{equation}
Trong công thức này, $\mathbf{x}$ là đặc trưng đầu vào, $\mathbf{y}$ là đầu ra của khối, và $\mathcal{F}$ biểu thị ánh xạ dư thừa cần được tối ưu hóa. Kết nối tắt giúp gradient lan truyền ngược dễ dàng qua các lớp mà không bị triệt tiêu (gradient vanishing). Phiên bản ResNet-18 có 18 lớp trọng số cấu tạo từ các khối BasicBlock (hai lớp tích chập $3\times3$ kết hợp với Batch Normalization và ReLU). Cấu trúc chi tiết của khối dư thừa được minh họa trong Hình~\ref{fig:resnet_arch}.

\begin{figure}[H]
    \centering
    \includegraphics[width=0.55\textwidth]{resnet_arch.png}
    \caption{Kiến trúc khối dư thừa (Residual Block) với kết nối tắt trong ResNet}
    \label{fig:resnet_arch}
\end{figure}

Dựa vào sơ đồ Hình~\ref{fig:resnet_arch}, đặc trưng đầu vào $\mathbf{x}$ đi qua hai nhánh song song. Nhánh thứ nhất thực hiện trích xuất thông qua hai lớp trọng số tích chập liên tiếp (kèm chuẩn hóa Batch Normalization và hàm kích hoạt ReLU). Nhánh thứ hai là kết nối tắt (Identity Shortcut) truyền trực tiếp $\mathbf{x}$ mà không chịu bất kỳ biến đổi phi tuyến tính nào. Hai nhánh sau đó được cộng chập trực tiếp trước khi đi qua lớp kích hoạt ReLU cuối cùng của khối. Trong trường hợp kích thước kênh hoặc độ phân giải không gian của $\mathbf{x}$ thay đổi giữa đầu vào và đầu ra, nhánh shortcut sẽ áp dụng một phép tích chập $1\times1$ để điều chỉnh số chiều tương thích. Thiết kế này giúp gradient trong quá trình lan truyền ngược có thể truyền thẳng qua nhánh shortcut mà không bị triệt tiêu bởi các phép nhân ma trận trọng số, cho phép mạng đạt độ sâu hàng trăm lớp một cách ổn định. Tổng số tham số của ResNet-18 là 11,18 triệu, cung cấp khả năng biểu diễn đặc trưng mạnh mẽ cho phổ đồ âm thanh.

\subsection{DenseNet-121}

Mạng tích chập kết nối dày đặc (DenseNet) do Huang và cộng sự đề xuất năm 2017 \cite{huang2017} đưa ra giải pháp kết nối trực tiếp tất cả các lớp với nhau trong một khối liên kết dày đặc (Dense Block). Thay vì thực hiện phép cộng đặc trưng như ResNet, mỗi lớp trong DenseBlock sẽ nhận bản đồ đặc trưng của toàn bộ các lớp phía trước làm đầu vào thông qua phép ghép nối (concatenation) theo chiều kênh:
\begin{equation}
    \mathbf{x}_l = H_l\left([\mathbf{x}_0, \mathbf{x}_1, \ldots, \mathbf{x}_{l-1}]\right)
    \label{eq:densenet}
\end{equation}
Trong đó, $[\mathbf{x}_0, \mathbf{x}_1, \ldots, \mathbf{x}_{l-1}]$ biểu diễn phép ghép bản đồ đặc trưng được tạo ra bởi các lớp trước đó, và $H_l$ là phép biến đổi phi tuyến tính. Cơ chế này giúp cải thiện đáng kể luồng thông tin và gradient trong mạng, khuyến khích tái sử dụng đặc trưng và giảm thiểu số tham số thực tế cần huấn luyện do các lớp không cần học lại đặc trưng trùng lặp. Giữa các khối Dense Block, mạng sử dụng các lớp chuyển tiếp (Transition Layers) gồm tích chập $1\times1$ và gộp trung bình $2\times2$. Sơ đồ khối liên kết dày đặc được biểu diễn trong Hình~\ref{fig:densenet_arch}.

\begin{figure}[H]
    \centering
    \includegraphics[width=0.75\textwidth]{densenet_arch.png}
    \caption{Sơ đồ liên kết dày đặc (Dense Block) trong kiến trúc DenseNet}
    \label{fig:densenet_arch}
\end{figure}

Phân tích sơ đồ Hình~\ref{fig:densenet_arch} cho thấy cấu trúc dòng chảy thông tin cực kỳ cấu trúc: Layer $H_1$ nhận đầu vào từ Layer $0$; Layer $H_2$ nhận đầu vào từ cả Layer $0$ và đặc trưng trích xuất của Layer $1$ sau khi ghép kênh; quá trình này tiếp diễn tuần tự cho đến Layer cuối cùng. Tham số growth rate ($k=32$) quy định số lượng kênh đặc trưng mới được tạo ra ở mỗi bước biến đổi. Thiết kế ghép kênh thay vì cộng chập giúp mô hình giữ lại toàn vẹn các thông tin tần số thấp và tần số cao được học từ các lớp trước đó mà không bị mờ nhòe đặc trưng, tối ưu hóa việc phân tách các đặc trưng âm học nhạy cảm. Kiến trúc DenseNet-121 gồm 121 lớp với tổng số 6,96 triệu tham số giúp khai thác triệt để các đặc trưng âm học ở nhiều cấp độ phân giải.

\subsection{EfficientNet-B0}

Mạng EfficientNet do Tan và Lê đề xuất năm 2019 \cite{tan2019} tối ưu hóa hiệu năng CNN một cách có hệ thống thông qua kỹ thuật tỷ lệ hóa hỗn hợp (Compound Scaling). Thay vì tăng riêng lẻ độ sâu, chiều rộng hoặc độ phân giải ảnh đầu vào, phương pháp này tỷ lệ hóa đồng đều cả ba chiều bằng một hệ số tỷ lệ $\phi$ cố định duy nhất:
\begin{equation}
    \text{depth}: d = \alpha^\phi, \quad \text{width}: w = \beta^\phi, \quad \text{resolution}: r = \gamma^\phi
    \label{eq:efficientnet}
\end{equation}
với ràng buộc $\alpha \cdot \beta^2 \cdot \gamma^2 \approx 2$ nhằm đảm bảo khi tăng $\phi$, lượng FLOPS yêu cầu tăng tuyến tính theo mức $2^\phi$. Khối MBConv (Mobile Inverted Bottleneck Convolution) là thành phần cấu tạo chính của EfficientNet, tích hợp tích chập tách biệt chiều sâu (depthwise separable convolution) kết hợp với khối chú ý Squeeze-and-Excitation (SE) giúp mô hình tự học trọng số biểu thị tầm quan trọng của từng kênh đặc trưng. Sơ đồ khối MBConv được minh họa trong Hình~\ref{fig:efficientnet_arch}.

\begin{figure}[H]
    \centering
    \includegraphics[width=0.8\textwidth]{efficientnet_arch.png}
    \caption{Kiến trúc khối MBConv tích hợp cơ chế chú ý Squeeze-and-Excitation trong EfficientNet}
    \label{fig:efficientnet_arch}
\end{figure}

Sơ đồ Hình~\ref{fig:efficientnet_arch} mô tả chi tiết luồng xử lý của khối MBConv. Đặc trưng đầu vào trước hết được mở rộng số kênh qua lớp tích chập $1\times1$ để chuyển sang không gian đặc trưng số chiều lớn, sau đó đi qua lớp tích chập tách biệt chiều sâu (Depthwise Conv) kích thước $3\times3$ hoặc $5\times5$ để học các thông tin không gian cục bộ một cách tiết kiệm chi phí tính toán. Đặc trưng tiếp tục đi qua khối Squeeze-and-Excitation (SE) thực hiện ép kênh qua Global Average Pooling và tái khôi phục qua hai lớp Fully Connected để tạo ra trọng số chú ý cho từng kênh, nhân trực tiếp vào đặc trưng nhằm làm nổi bật các kênh tần số âm thanh quan trọng. Cuối cùng, lớp chiếu $1\times1$ hạ số chiều kênh về mức ban đầu và được cộng chập với shortcut từ đầu vào. Thiết kế MBConv giúp EfficientNet-B0 (4,01 triệu tham số) đạt hiệu năng cao nhất với chi phí tính toán tối thiểu trên phổ đồ giọng nói.

\subsection{MobileNet-V2}

MobileNet-V2 được Sandler và cộng sự giới thiệu vào năm 2018 \cite{sandler2018} là một kiến trúc CNN gọn nhẹ được tối ưu hóa cho các ứng dụng chạy trên thiết bị di động và hệ thống nhúng. Mô hình sử dụng tích chập tách biệt theo chiều sâu (depthwise separable convolution) để phân tách phép tích chập thông thường thành hai bước tích chập độc lập gồm depthwise convolution và pointwise convolution, giúp giảm mạnh số lượng phép tính. MobileNet-V2 giới thiệu khối nghịch đảo residual (Inverted Residual Block) thực hiện mở rộng số kênh ở lớp giữa và chiếu giảm số kênh ở đầu ra thông qua lớp cổ chai tuyến tính (Linear Bottleneck) không sử dụng hàm kích hoạt phi tuyến tính nhằm tránh mất mát thông tin hữu ích trong không gian số chiều thấp. Cấu trúc khối Inverted Residual được thể hiện trong Hình~\ref{fig:mobilenet_arch}.

\begin{figure}[H]
    \centering
    \includegraphics[width=0.8\textwidth]{mobilenet_arch.png}
    \caption{Kiến trúc khối nghịch đảo residual (Inverted Residual Block) và Linear Bottleneck trong MobileNet-V2}
    \label{fig:mobilenet_arch}
\end{figure}

Phân tích Hình~\ref{fig:mobilenet_arch} làm rõ sự khác biệt giữa Inverted Residual và Residual Block thông thường. Thay vì kết nối các lớp có số lượng kênh lớn, nhánh shortcut của MobileNet-V2 kết nối trực tiếp giữa các lớp cổ chai tuyến tính (Linear Bottleneck) có số lượng kênh rất nhỏ (không gian số chiều thấp). Trong khối, đặc trưng trước tiên được mở rộng kênh lên gấp 6 lần (Expansion Layer), sau đó trích xuất không gian bằng tích chập tách biệt chiều sâu (Depthwise Convolution), và cuối cùng chiếu giảm số kênh (Projection Layer) sử dụng hàm kích hoạt tuyến tính (Linear Activation) thay vì ReLU. Thiết kế Linear Bottleneck cực kỳ quan trọng vì hàm ReLU trong không gian số chiều thấp có xu hướng triệt tiêu hoàn toàn các giá trị âm, làm mất mát nghiêm trọng thông tin đặc trưng; việc giữ kích hoạt tuyến tính ở đầu ra lớp cổ chai giúp duy trì nguyên vẹn các đặc trưng âm học nhạy cảm. Với chỉ 2,23 triệu tham số, MobileNet-V2 đáp ứng khả năng tính toán thời gian thực trên các phần cứng cấu hình yếu.


\section{Hàm lỗi Focal Loss}
\label{sec:focal_loss}

Hàm lỗi Cross Entropy truyền thống xử lý mọi mẫu dữ liệu một cách đồng đều:
\begin{equation}
    \text{CE}(p_t) = -\log(p_t)
    \label{eq:ce}
\end{equation}
Điều này dẫn đến việc mô hình dễ bị chi phối bởi lớp đa số (Neutral) trong môi trường dữ liệu mất cân bằng nghiêm trọng. Nhằm khắc phục hạn chế này, Focal Loss \cite{lin2017} giới thiệu thêm một hệ số điều chỉnh động $(1 - p_t)^\gamma$ vào công thức tính lỗi:
\begin{equation}
    \text{FL}(p_t) = -\alpha_t \cdot (1 - p_t)^\gamma \cdot \log(p_t)
    \label{eq:focal}
\end{equation}
Trong công thức trên, $p_t$ đại diện cho xác suất dự đoán đúng của mô hình đối với lớp thực tế; $\gamma \geq 0$ là hệ số tập trung (focusing parameter) dùng để điều khiển tốc độ giảm trọng số lỗi đối với các mẫu dễ phân loại (trong đồ án sử dụng giá trị tối ưu $\gamma = 2.0$); và $\alpha_t$ là hệ số cân bằng lớp. Cần lưu ý rằng trong bài báo gốc của Lin và cộng sự \cite{lin2017}, Focal Loss được thiết kế ban đầu cho bài toán phân loại nhị phân với $\alpha \in [0, 1]$ là một siêu tham số cố định được điều chỉnh thủ công. Để phù hợp với bài toán phân loại đa lớp thực tế trong đồ án này, hệ số $\alpha_t$ được triển khai dưới dạng một cơ chế gán trọng số lớp động (class weighting scheme) dựa trên nghịch đảo tần suất của từng lớp trong tập huấn luyện:
\begin{equation}
    \alpha_c = \frac{N_{\text{total}}}{C \times N_c}
\end{equation}
trong đó $N_c$ là số mẫu của lớp $c$, $C$ là tổng số lớp ($C = 4$), và $N_{\text{total}}$ là tổng số mẫu của tất cả các lớp trong tập huấn luyện. Cơ chế điều chỉnh động này giúp giảm mạnh phần đóng góp lỗi từ các mẫu đã phân loại đúng và buộc mô hình tập trung tối ưu hóa trên các mẫu khó thuộc lớp thiểu số (như Happiness, Anger).


\section{Kỹ thuật tăng cường dữ liệu EMix}
\label{sec:emix}

EMix (Emotion-aware Mixup) \cite{dang2023emix} là kỹ thuật tăng cường dữ liệu dựa trên nền tảng của thuật toán MixUp truyền thống \cite{zhang2018}, được thiết kế chuyên biệt cho bài toán nhận dạng cảm xúc giọng nói (SER). Thay vì thực hiện trộn ngẫu nhiên hai mẫu bất kỳ tạo ra các nhãn mập mờ làm mờ ranh giới cảm xúc, EMix đề xuất các chiến lược trộn có kiểm soát dựa trên nhãn cảm xúc thực tế.

Chiến lược thứ nhất là EMix-S (Same-class Mixup), thực hiện trộn tuyến tính hai mẫu âm thanh có cùng nhãn cảm xúc để củng cố ranh giới và tính phân tách của từng lớp. Công thức toán học trộn được biểu diễn:
\begin{equation}
    \tilde{x}_i = \lambda \cdot x_i + (1 - \lambda) \cdot x_j, \quad \text{với } y_i = y_j, \quad \lambda \sim U(0, 1)
    \label{eq:emix_s}
\end{equation}
Trong đó, hai phổ đồ $x_i$ và $x_j$ có cùng nhãn cảm xúc ($y_i = y_j$) được trộn với nhau thông qua trọng số $\lambda$ được lấy ngẫu nhiên từ phân phối đều $U(0, 1)$, và nhãn đầu ra cho mẫu mới được giữ nguyên là $\tilde{y}_i = y_i$.

Chiến lược thứ hai là EMix-N (Neutral Mixup), thực hiện trộn mẫu cảm xúc mục tiêu với một mẫu cảm xúc trung tính (Neutral) đóng vai trò như một dạng nhiễu nền tự nhiên. Phép trộn này giúp mô hình nhận diện tốt hơn các đặc trưng cảm xúc ngay cả khi bị suy hao hoặc bị lấn át bởi trạng thái nói bình thường. Công thức cụ thể là:
\begin{equation}
    \tilde{x}_i = \lambda \cdot x_i + (1 - \lambda) \cdot x_{\text{neu}}, \quad \lambda \sim U(0.5, 1.0)
    \label{eq:emix_n}
\end{equation}
Với hệ số $\lambda \geq 0.5$ nhằm đảm bảo đặc trưng cảm xúc mục tiêu vẫn chiếm ưu thế ngữ nghĩa so với mẫu Neutral, và nhãn đích đầu ra vẫn được duy trì theo mẫu cảm xúc mục tiêu.

Chiến lược cuối cùng là EMix-NS (Combined) đại diện cho phương pháp kết hợp đồng thời cả hai chiến lược trên. Với mỗi mini-batch trong quá trình huấn luyện, thuật toán sẽ ngẫu nhiên lựa chọn thực hiện EMix-N hoặc EMix-S với xác suất đồng đều ($p = 0.5$). Sự kết hợp này mang lại dải biến thể phong phú cho tập dữ liệu huấn luyện, giúp mô hình cải thiện tối đa tính tổng quát.
